% ECONOMICS_PROBLEM: {"id": "I22-533", "title": "Effective uses of school funding", "jel_code": "I22", "source_id": "OP-0452"}
% !TeX program = xelatex
\documentclass[11pt,a4paper]{article}
\usepackage[margin=23mm]{geometry}
\usepackage{fontspec}
\setmainfont{Latin Modern Roman}
\usepackage{amsmath,amssymb,mathtools}
\usepackage{xcolor,enumitem,booktabs,longtable,array,xurl,hyperref}
\definecolor{muted}{HTML}{53636C}
\hypersetup{colorlinks=true,urlcolor=blue,linkcolor=blue}
\setlength{\parindent}{0pt}
\setlength{\parskip}{5pt}
\newcommand{\R}{\mathbb R}
\newcommand{\E}{\mathbb E}
\newcommand{\N}{\mathbb N}
\newcommand{\ind}{\mathbf 1}
\newcommand{\Var}{\operatorname{Var}}
\newcommand{\Cov}{\operatorname{Cov}}
\newcommand{\supp}{\operatorname{supp}}
\DeclareMathOperator*{\argmin}{arg\,min}
\DeclareMathOperator*{\argmax}{arg\,max}
\newcommand{\entrymeta}[3]{{\sffamily\small\color{muted}\textbf{\texttt{#1}}\quad|\quad #2\par #3\par}}
\newcommand{\Problemref}[1]{\texttt{#1}}
\newcommand{\JELref}[2]{\texttt{#2} (JEL \texttt{#1})}
\begin{document}
\section*{Effective uses of school funding}
\textbf{Problem ID:} \texttt{I22-533}\quad \textbf{JEL:} \texttt{I22}\par
% BEGIN ECONOMICS STATEMENT
\label{problem:OP-0452}
{\sffamily\small\color{muted}Original subject group: Education, families and demography\par}
\label{definitions:problem:30007734}
\textbf{Audit classification:} Published research agenda. Checked source identity and appropriate agenda/background support; expanded intervention, units, population, definedness and causal restrictions. Exact targets and variants are editorial.
\subsection*{Definitions and precise research target}
\textit{Editorial formalization.} Consider public schools serving disadvantaged pupils in one state over five budget years. Let \(b>0\) be additional annual funding per pupil and \(x=(x_1,\ldots,x_J)\) a spending allocation across \(J\) prespecified uses, such as tutoring, teacher retention, class-size reduction, and facilities. Feasibility requires \(x_j\geq0\) and \(\sum_{j=1}^{J}x_j=b\). Let \(Y_i(x)\) denote pupil \(i\)'s educational attainment and independently scaled learning outcomes, combined using fixed weights announced before estimation. Define
\[x^*(b)\in\arg\max_{x\geq0:\sum_jx_j=b}E[Y_i(x)-Y_i(0)].\]
Assume finite expected outcomes and a continuous mean response on the compact allocation simplex, or report a supremum if attainment cannot be justified. Estimate credible response surfaces or feasible policy rankings using randomized spending packages or justified administrative variation. School context, staff availability, implementation quality, and baseline resources must be measured; average funding effects do not identify this allocation problem. Address substitution of local funds and spillovers from teacher recruitment. Report whether rankings transport to a separately specified district population, with overlap conditions and sensitivity to unobserved effect modifiers.

The allocation vector \(x\in\mathbb R_+^J\) is annual incremental real spending per baseline pupil repeated for five years; \(x=0\) is the original budget. Feasible uses require declared staffing, procurement, and capacity conditions; if some points of the simplex cannot be implemented, optimize over a stated closed subset \(\mathcal X_b\), assumed nonempty. A scalar outcome index is \(Y_i(x)=\alpha_1\operatorname{learning}_i(x)+\alpha_2\operatorname{attainment}_i(x)\), with fixed scales and nonnegative announced weights. Continuous finite mean responses on compact \(\mathcal X_b\) attain a maximum. This is routine optimization after identifying the response surface. Trial packages identify only their own contrasts; interpolation between untested uses is an additional response restriction. Fiscal substitution, teacher-market spillovers and eligible-cohort movement must be included or excluded explicitly. All expectations use a stated baseline population and probability law. Identification means every admissible data-generating model with the same observed-data law yields the same displayed target; otherwise report the identified set. Exact equations, horizons and assumptions are editorial, rather than source-given canonical conjectures.
\subsection*{Variants and assumption changes}
\begin{enumerate}
\item \textbf{Editorial integer-capacity variant.} Restrict spending to indivisible staffing and classroom packages. Optimize a finite feasible menu rather than an unrealistic full simplex.
\item \textbf{Editorial transport variant.} Integrate conditional responses against target law \(P_T(X)\) under overlap and conditional invariance. Bound unrepresented contexts when these assumptions fail.
\end{enumerate}
\subsection*{References and source audit}
\textbf{1.} Explicit difficulty generalizing spending effects without understanding contexts and uses; allocation optimizer is editorial. Locator: Evaluation Review 48(3), 461–494; author-page abstract. Full text or primary publication page inspected. URL: \url{https://www.hanushek.net/publications/contexts-convenience-generalizing-published-evaluations-school-finance-policies-1}.
\begin{enumerate}\raggedright
\item\hypertarget{definitions:source:30007734:1}{} Danielle V. Handel; Eric A. Hanushek. 2024. \emph{Contexts of Convenience: Generalizing from Published Evaluations of School Finance Policies}. Evaluation Review 48(3), 461–494; author-page abstract.
\url{https://www.hanushek.net/publications/contexts-convenience-generalizing-published-evaluations-school-finance-policies-1}.
DOI: \url{https://doi.org/10.1177/0193841X241228335}.
\end{enumerate}

\subsection*{Common conventions: Source mathematical conventions}

These conventions apply to each entry unless explicitly replaced.

\textbf{Models and identification.} A primitive model \(m\) specifies all probability laws, feasible choices, information, equilibrium and measurement rules used by its target. The admissible class \(\mathcal M\) is fixed independently of outcomes. If \(P_O(m)\) is its observable probability law and \(\tau(m)\) the target, its identified set at observable law \(P\) is
\[
 \mathcal I_\tau(P)=\{\tau(m):m\in\mathcal M,\ P_O(m)=P\}.
\]
Point identification means that this set is a singleton. An empty set signals incompatibility of data and assumptions. Coordinate bounds need not describe the joint set. Bounds are sharp only if every retained target is attainable or arbitrarily approachable by compatible models. Finite bounds require explicit restrictions on unobserved quantities; unrestricted confounding, hidden wealth, extrapolation or arbitrary equilibrium selection can yield uninformative sets.

\textbf{Causal interventions.} A potential outcome \(Y_i(p)\) is the outcome under a complete policy regime \(p\), including timing, financing, information and market spillovers. The target population and units are fixed before treatment. With interference, use the complete assignment vector or a justified exposure mapping; individual no-interference notation alone cannot identify an economy-wide intervention. A difference of mean potential outcomes does not require the cross-world joint distribution. Conditioning on post-treatment migration, survival, enrollment or employment changes the estimand and needs separate justification.

\textbf{Statistical statements.} An estimator, confidence set, forecast or test specifies a sampling law, model class and loss/coverage criterion. Uniform \((1-\alpha)\)-coverage means \(\inf_{m\in\mathcal M}\Pr_m\{\tau(m)\in C_n\}\geq1-\alpha\), or an explicitly asymptotic version. The whole parameter space trivially has coverage, so informativeness requires a stated width, risk or power objective. A forecasting improvement is relative to a named benchmark, information set and evaluation population.

\textbf{Equilibrium and optimization.} An equilibrium comprises feasible plans, optimality, beliefs where needed, budget constraints and all market/resource clearing. The primitives or a stated benchmark define each correspondence \(\mathcal E(p;m)\). Nonemptiness is assumed or proved, never inferred just from compact policy sets. A selected-equilibrium target uses a declared selection rule; a robust target quantifies over all permitted equilibria. Use suprema or infima unless attainment is established. An \(\varepsilon\)-optimal policy is feasible and within \(\varepsilon\) of the finite optimum value. Report an empty feasible set separately.

\textbf{Welfare and accounting.} Normative utility, interpersonal normalization, social weights, numeraire, discounting and the use of public receipts are primitives. Behavioral utility need not coincide with the welfare criterion. Household welfare includes ownership income and rents once; adding profits again to compensating variation generally double-counts them. A monetary equivalent is defined by an explicit transfer and price convention, with monotonicity and range conditions. Average effects, mechanism contributions, transfers and welfare losses are different targets. Infinite sums require convergence and dynamic feasible plans require terminal or no-Ponzi conditions.

\textbf{Source versions.} A working-paper date, revision date and journal publication year are distinguished. A DOI for a journal version is not automatically the DOI of an earlier manuscript. Locators refer to the stated version and printed pages where specified. A metadata-only check does not verify a claim about a full-text conclusion. Every entry's source subsection narrows the provenance claim accordingly.
% END ECONOMICS STATEMENT
\end{document}
